% ============================================================
%  SECCIÓN — PLANO Y DETALLE DE CIMENTACIÓN (ARQUETA)
% ============================================================
\section{Plano de cimentación y perfil del terreno}
\label{sec:arqueta_cimentacion}

% -----------------------------------------------------------
\subsection{Plano acotado --- sección transversal}
\label{subsec:arqueta_seccion}
% -----------------------------------------------------------

La figura~\ref{fig:arqueta_seccion} muestra la sección transversal de la
arqueta de fangos mezcla con todas las dimensiones adoptadas y el perfil
estratigráfico del Sondeo~1.

\begin{figure}[H]
\centering
\begin{tikzpicture}[
  scale=1.80, >=Stealth,
  dim/.style    = {<->, thin, NavyBlue},
  dlbl/.style   = {font=\scriptsize, midway, fill=white, inner sep=1pt},
  soil1a/.style = {fill=Goldenrod!20},
  soil1b/.style = {fill=Goldenrod!45},
  soil2/.style  = {fill=BurntOrange!18},
  concrete/.style = {fill=gray!22},
]

%% ---- RIGHT STRATIGRAPHY ----
\fill[soil1a] ( 1.9, 0.0) rectangle ( 4.2,-2.5);
\fill[soil1b] ( 1.9,-2.5) rectangle ( 4.2,-3.5);
\fill[soil2]  ( 1.9,-3.5) rectangle ( 4.2,-4.2);
\draw[thin, gray!60] ( 1.9,-2.5) -- ( 4.2,-2.5);
\draw[thin, gray!60] ( 1.9,-3.5) -- ( 4.2,-3.5);
\draw[thick] ( 1.9, 0.0) -- ( 4.2, 0.0);
\fill[pattern=north east lines, pattern color=gray!40]
  ( 1.9, 0.0) rectangle ( 4.2, 0.25);

%% ---- LEFT SOIL ----
\fill[soil1a] (-3.8, 0.0) rectangle (-1.9,-2.5);
\fill[soil1b] (-3.8,-2.5) rectangle (-1.9,-3.5);
\fill[soil2]  (-3.8,-3.5) rectangle (-1.9,-4.2);
\draw[thin, gray!60] (-3.8,-2.5) -- (-1.9,-2.5);
\draw[thin, gray!60] (-3.8,-3.5) -- (-1.9,-3.5);
\draw[thick] (-3.8, 0.0) -- (-1.9, 0.0);
\fill[pattern=north east lines, pattern color=gray!40]
  (-3.8, 0.0) rectangle (-1.9, 0.25);

%% ---- RASANTE (dashed over arqueta gap) ----
\draw[gray!40, dashed] (-1.9, 0.0) -- ( 1.9, 0.0);

%% ---- ARQUETA CONCRETE (losa + paredes + tapa) ----
\fill[concrete] (-1.9,-1.8) rectangle ( 1.9, 2.3);
%% Interior void (paredes huecas, tapa cierra)
\fill[white]    (-1.5,-1.5) rectangle ( 1.5, 2.0);

%% Outlines
\draw[thick] (-1.9,-1.8) -- (-1.9, 2.3);   % left outer
\draw[thick] (-1.5,-1.5) -- (-1.5, 2.0);   % left inner wall
\draw[thick] ( 1.9,-1.8) -- ( 1.9, 2.3);   % right outer
\draw[thick] ( 1.5,-1.5) -- ( 1.5, 2.0);   % right inner wall
\draw[thick] (-1.9,-1.8) -- ( 1.9,-1.8);   % bottom losa (intradós)
\draw[thick] (-1.9,-1.5) -- (-1.5,-1.5);   % top losa left (trasdós)
\draw[thick] ( 1.5,-1.5) -- ( 1.9,-1.5);   % top losa right (trasdós)
\draw[thick] (-1.9, 2.0) -- ( 1.9, 2.0);   % bottom of tapa
\draw[thick] (-1.9, 2.3) -- ( 1.9, 2.3);   % coronación (top of tapa)

%% ---- NIVEL FREATICO ----
\draw[NavyBlue!80, semithick, dashed] (-3.8,-2.5) -- ( 4.2,-2.5);
\node[NavyBlue!80, right, font=\scriptsize\bfseries]
  at (4.2,-2.5) {N.F.\ $= -2{,}50$~m};

%% ---- DIMENSION LINES ----

%% D1: Interior width B_int = 3.00 m
\draw[dim] (-1.5,-0.80) -- (1.5,-0.80)
  node[dlbl, above] {$B_{\mathrm{int}} = 3{,}00$~m};
\draw[thin,gray!50] (-1.5,-0.60) -- (-1.5,-0.90);
\draw[thin,gray!50] ( 1.5,-0.60) -- ( 1.5,-0.90);

%% D2: Exterior width B_ext = 3.60 m
\draw[dim] (-1.9,-2.25) -- (1.9,-2.25)
  node[dlbl, above] {$B_{\mathrm{ext}} = 3{,}60$~m};
\draw[thin,gray!50] (-1.9,-2.00) -- (-1.9,-2.35);
\draw[thin,gray!50] ( 1.9,-2.00) -- ( 1.9,-2.35);

%% D3+D4: Wall thickness = 0.30 m (both sides)
\draw[dim] (-1.9, 2.55) -- (-1.5, 2.55)
  node[dlbl, above] {$e\!=\!0{,}30$};
\draw[dim] ( 1.5, 2.55) -- ( 1.9, 2.55)
  node[dlbl, above] {$e\!=\!0{,}30$};

%% D5: Interior height H_int = 3.50 m (right side)
\draw[dim] (2.15,-1.5) -- (2.15, 2.0)
  node[dlbl, right=2pt] {$H_{\mathrm{int}} = 3{,}50$~m};
\draw[thin,gray!50] (1.95,-1.5) -- (2.30,-1.5);
\draw[thin,gray!50] (1.95, 2.0) -- (2.30, 2.0);

%% D6: Losa thickness e_L = 0.30 m (right side)
\draw[dim] (2.15,-1.8) -- (2.15,-1.5)
  node[dlbl, right=2pt] {$e_L\!=\!0{,}30$};
\draw[thin,gray!50] (1.95,-1.8) -- (2.30,-1.8);

%% D6b: Tapa thickness = 0.30 m (right side, above D5)
\draw[dim] (2.15, 2.0) -- (2.15, 2.3)
  node[dlbl, right=2pt] {$e_T\!=\!0{,}30$};
\draw[thin,gray!50] (1.95, 2.3) -- (2.30, 2.3);

%% D7: Excavation depth to trasdós = 1.50 m (left side)
\draw[dim] (-2.50, 0.0) -- (-2.50,-1.5)
  node[dlbl, left=2pt] {$1{,}50$~m};
\draw[thin,gray!50] (-2.65, 0.0) -- (-2.35, 0.0);
\draw[thin,gray!50] (-2.65,-1.5) -- (-2.35,-1.5);

%% D8: Df = 1.80 m (left side, rasante to intradós)
\draw[dim] (-3.15, 0.0) -- (-3.15,-1.8)
  node[dlbl, left=2pt] {$D_f = 1{,}80$~m};
\draw[thin,gray!50] (-3.30, 0.0) -- (-3.00, 0.0);
\draw[thin,gray!50] (-3.30,-1.8) -- (-3.00,-1.8);

%% D9: Height above grade to coronation = 2.30 m (left side)
\draw[dim] (-2.10, 0.0) -- (-2.10, 2.3)
  node[dlbl, left=2pt] {$2{,}30$~m};

%% ---- STRATIGRAPHY LABELS ----
\node[font=\scriptsize, align=left, text width=1.8cm]
  at (3.05,-0.80) {\textbf{E-1a}\\ Arena limosa\\ no saturada\\ $\phi'=29°$};
\node[font=\scriptsize, align=left, text width=1.8cm]
  at (3.05,-3.00) {\textbf{E-1b}\\ Arena limosa\\ saturada};
\node[font=\scriptsize, align=left, text width=1.8cm]
  at (3.05,-3.85) {\textbf{E-2}\\ Arcillas\\ versicolores};

%% ---- INTERIOR LABELS ----
\node[gray!80, font=\footnotesize, align=center] at (0.0, 0.7)
  {$L\times B\times H = 4{,}00\times3{,}00\times3{,}50$~m};
\node[gray!80, font=\scriptsize] at (0.0, 0.2)
  {HA-30/B/20/IIb+Qc --- cem.\ SR};

%% Tapa label
\node[gray!80, font=\scriptsize] at (0.0, 2.15)
  {Tapadera $e=0{,}30$~m};

%% ---- ELEVATION REFERENCE LABELS ----
\node[gray!70, font=\tiny, left] at (-3.8, 2.3)  {$+2{,}30$~m};
\node[gray!70, font=\tiny, left] at (-3.8, 0.0)  {$\pm0{,}00$};
\node[gray!70, font=\tiny, left] at (-3.8,-1.5)  {$-1{,}50$~m};
\node[gray!70, font=\tiny, left] at (-3.8,-1.8)  {$-1{,}80$~m};
\node[gray!70, font=\tiny, left] at (-3.8,-2.5)  {$-2{,}50$~m};

\end{tikzpicture}
\caption{Sección transversal de la arqueta de fangos mezcla con perfil
         estratigráfico del Sondeo~1. Dimensiones en metros.
         $D_f=1{,}80$~m (intradós); trasdós superior losa: $-1{,}50$~m;
         coronación: $+2{,}30$~m; N.F.: $-2{,}50$~m.}
\label{fig:arqueta_seccion}
\end{figure}

% -----------------------------------------------------------
\subsection{Justificación de la cota de cimentación}
\label{subsec:cota_cim}
% -----------------------------------------------------------

La cota de apoyo de la losa se fija en $D_f=1{,}80$~m (intradós) según la
convención del CTE DB SE-C, en virtud de los tres condicionantes siguientes:

\subsubsection*{(a) Estrato portante}

El Estrato~E-1a (arena limosa, $\phi'=29°$, $c'=0$) ofrece una capacidad
portante $q_u=101{,}77$~Tn/m² y un factor de seguridad $FS=15{,}2\gg3{,}0$
(sección~\ref{subsec:cap_portante_arqueta}).
La profundización hasta E-2 (arcillas, $D_f\approx4{,}0$~m) no aporta
mejora estructural y requiere atravesar el nivel freático,
con el consiguiente coste de agotamiento y entibación.

\subsubsection*{(b) Nivel freático}

El N.F.\ se localiza a $-2{,}50$~m, $0{,}70$~m por debajo del intradós
($-1{,}80$~m). La cimentación se ejecuta completamente en \textbf{seco},
sin necesidad de agotamiento permanente.
En la hipótesis pésima de N.F.\ en rasante, el factor de seguridad a la
flotación es $FS_\text{flot}=2{,}17>1{,}05$~\checkmark
(sección~\ref{subsec:flotacion}).

\subsubsection*{(c) Condicionante medioambiental}

Al estar en zona medioambientalmente protegida se minimiza el volumen de
excavación (base de la losa a tan solo $1{,}80$~m de profundidad) y se
evita el empleo de lodos bentoníticos, incompatibles con la zona protegida.
La estanqueidad se garantiza con el sistema de impermeabilización descrito
en la sección~\ref{subsec:impermeabilizacion_arqueta}.

\begin{supuestobox}[Cota de cimentación adoptada --- Arqueta]
  \begin{itemize}
    \item Cota intradós (base de la losa): $-1{,}80$~m $\Rightarrow$
          \textbf{$D_f=1{,}80$~m} (CTE DB SE-C).
    \item Cota trasdós superior de la losa: $-1{,}50$~m.
    \item Exceso sobre N.F.: $\Delta h = 2{,}50-1{,}80 = \mathbf{0{,}70}$~m
          (cimentación en seco)~\checkmark.
    \item Capa de HL-150 de limpieza: $e=10$~cm
          (cota fondo de excavación $=-1{,}90$~m).
  \end{itemize}
\end{supuestobox}

% -----------------------------------------------------------
\subsection{Medidas frente a la agresividad química}
\label{subsec:impermeabilizacion_arqueta}
% -----------------------------------------------------------

\begin{notabox}[Protección frente a sulfatos y filtración de fangos]
  Concentración de $3\,500$~mg~SO$_4^{2-}$/kg de suelo seco
  $\rightarrow$ clase de exposición \textbf{Qc} (EHE-08).
  Se adoptarán las siguientes medidas:
  \begin{enumerate}
    \item \textbf{Hormigón con cemento resistente a sulfatos.}
          HA-30/B/20/IIb+Qc con \textbf{cemento SR} (CEM~I~42,5~SR),
          relación $a/c\leq0{,}50$, cemento mínimo $\geq325$~kg/m³,
          recubrimiento nominal 45~mm.
    \item \textbf{Impermeabilización exterior} de losa y paredes enterradas
          mediante lámina HDPE ($e\geq2$~mm), sobre geotextil de
          protección no tejido ($150$~g/m²).
    \item \textbf{Lámina drenante nodular} (tipo DELTA-MS o similar) sobre
          la HDPE para conducir el agua infiltrada al tubo perimetral.
    \item \textbf{Tubo dren perimetral} Ø~100~mm ranurado, rodeado de grava
          lavada $\phi~20/40$~mm filtrada con geotextil, en la base del
          trasdós de las paredes enterradas.
    \item \textbf{Relleno de trasdós} con zahorra artificial compactada en
          tongadas de 30~cm hasta rasante.
  \end{enumerate}
  El sistema garantiza la \textbf{doble barrera} (impermeabilización + drenaje)
  exigida por la clase Qc y la normativa medioambiental.
\end{notabox}

% -----------------------------------------------------------
\subsection{Prescripciones técnicas adicionales}
\label{subsec:prescripciones_arqueta}
% -----------------------------------------------------------

\begin{itemize}
  \item \textbf{Excavación.} Se prevé entibación o talud de seguridad
        ($1\text{H}:1\text{V}$ mínimo) durante la excavación hasta
        $-1{,}90$~m (incluyendo hormigón de limpieza).
  \item \textbf{Control geotécnico.} Se recomienda confirmación in~situ
        de los parámetros de E-1a antes del hormigonado de la losa
        (ensayo de carga de placa o densímetro nuclear).
  \item \textbf{Hormigón de limpieza.} Se hormigonará sobre $10$~cm de
        HL-150 para evitar la contaminación del hormigón estructural
        con finos del terreno.
  \item \textbf{Juntas de hormigonado.} Las juntas entre losa y paredes
        incorporarán banda de PVC tipo \emph{waterstop} para garantizar
        la estanqueidad ante el ascenso eventual del N.F.
\end{itemize}

% -----------------------------------------------------------
\subsection{Resumen de la solución de cimentación}
\label{subsec:resumen_cim_arqueta}
% -----------------------------------------------------------

\begin{table}[H]
  \centering
  \caption{Resumen de la cimentación adoptada para la arqueta de fangos mezcla}
  \label{tab:resumen_cim_arqueta}
  \begin{tabular}{@{}lll@{}}
    \toprule
    \textbf{Parámetro} & \textbf{Valor} & \textbf{Verificación}\\
    \midrule
    Tipo de cimentación    & Losa de H.A.\ superficial   & ---\\
    Dimensiones losa       & $4{,}60\times3{,}60$~m      & Dimensiones exteriores\\
    Espesor losa           & $e=0{,}30$~m                & Igual que paredes/tapa\\
    $D_f$ (intradós)       & $1{,}80$~m                  & Excav.\ en seco \checkmark\\
    Trasdós superior       & $-1{,}50$~m                 & $D_f-e_L$\\
    Estrato de apoyo       & E-1a (arena limosa)         & $\phi'=29°$, $c'=0$\\
    \midrule
    $\sigma_\text{contacto}$ & $6{,}70$~Tn/m²  & ---\\
    Cap.\ portante $q_u$     & $101{,}77$~Tn/m² & ---\\
    $FS$ hundimiento         & $15{,}2\geq3{,}0$ & \textbf{OK}~\checkmark\\
    \midrule
    Presión neta $\sigma_\text{neta}$ & $3{,}42$~Tn/m²  & compens.\ 49\%\\
    Asiento total estimado  & $s=2{,}51$~cm     & $\leq5{,}0$~cm \textbf{OK}~\checkmark\\
    $FS$ flotación (pésimo) & $2{,}17\geq1{,}05$ & \textbf{OK}~\checkmark\\
    \midrule
    Hormigón               & HA-30/B/20/IIb+Qc  & Cem.\ SR\\
    Acero                  & B~500~SD            & ---\\
    Recubrimiento nominal  & 45~mm               & Exp.\ Qc\\
    Impermeabilización     & HDPE $e\geq2$~mm    & + drenante + dren Ø100\\
    \bottomrule
  \end{tabular}
\end{table}
